% GENERATED FILE. Edit canonical lessons, then run npm run generate:books.
% canonical-source-hash: fnv1a64:7234ab04a3b0ecae
% canonical-lessons: JA-C197-shiyakusho, JA-C197-hiroba, JA-C197-koujou, JA-C197-tateru, JA-C197-mitsukaru, JA-R197-review-196-197, JA-R197-review-final-197

\chapter{City Hall, Squares and Factories}
\label{ch:ja-city-hall-squares-and-factories}
\hlchaptermodality{197}

\begin{chapteropening}
\textbf{By the end of this chapter:} \emph{I can talk about public places and new buildings in town.}

Complete the last lesson of chapter 197: I can talk about public places and new buildings in town.
\end{chapteropening}

\section[\texorpdfstring{\ja{しやくしょ} (shiyakusho)}{しやくしょ (shiyakusho)}]{\ja{しやくしょ} (shiyakusho) — city hall; municipal office}

\label{lesson:JA-C197-shiyakusho}



\begin{quote}
Before the new one: say the Japanese for map, then the Japanese for to go round.
\end{quote}

\subsection*{You'll want to know: \ja{しやくしょ}}
\textbf{\ja{しやくしょ}} — \emph{shiyakusho} — \textquotedblleft{}city hall; municipal office\textquotedblright{}.

\begin{quote}
\textbf{\ja{しやくしょは} \ja{えきの} \ja{みなみに} \ja{あります。}} — \emph{shiyakusho wa eki no minami ni arimasu} — \textquotedblleft{}City hall is south of the station.\textquotedblright{}
\end{quote}

\begin{cousinweb}
\textbf{\ja{しやくしょ}} (\emph{shiyakusho}) is Sino-Japanese (kango): three Chinese characters read with their on'yomi, readings borrowed from Chinese. They are \emph{shi}, \textquotedblleft{}market, city\textquotedblright{} (Middle Chinese \emph{dzyiX}), \emph{yaku}, \textquotedblleft{}duty, office\textquotedblright{} (\emph{ywek}) and \emph{sho}, \textquotedblleft{}place\textquotedblright{} (\emph{srjoX}) (KANJIDIC2; Middle Chinese in Baxter's transcription).
\end{cousinweb}

\begin{grammarlens}[title={What you've built}]
One more word for everyday life at A2.
\end{grammarlens}

\subsection*{Your turn}
\begin{itemize}
\raggedright
  \item \emph{Say it:} \emph{shiyakusho}
  \item \emph{Say it:} \emph{shiyakusho}, once more
  \item \emph{Say it:} \emph{shiyakusho wa eki no minami ni arimasu} — City hall is south of the station.
  \item \emph{Recall:} say the Japanese for map, then the Japanese for to go round, then say \emph{shiyakusho} again
\end{itemize}

\begin{tcolorbox}[breakable,colback=teal!4,colframe=teal!35!black,title={Before you move on}]
What does \ja{しやくしょ} mean? (\textquotedblleft{}City hall; municipal office\textquotedblright{}.) Say it once more.
\end{tcolorbox}
\section[\texorpdfstring{\ja{ひろば} (hiroba)}{ひろば (hiroba)}]{\ja{ひろば} (hiroba) — town square; open space, plaza}

\label{lesson:JA-C197-hiroba}



\begin{quote}
Before the new one: say the Japanese for to go round, then the Japanese for city hall.
\end{quote}

\subsection*{You'll want to know: \ja{ひろば}}
\textbf{\ja{ひろば}} — \emph{hiroba} — \textquotedblleft{}town square; open space, plaza\textquotedblright{}.

\begin{quote}
\textbf{\ja{ひろばで} \ja{まつりが} \ja{あります。}} — \emph{hiroba de matsuri ga arimasu} — \textquotedblleft{}There is a festival in the square.\textquotedblright{}
\end{quote}

\begin{cousinweb}
\textbf{\ja{ひろば}} (\emph{hiroba}) is native Japanese, built from \emph{hiro}, the stem of \emph{hiroi}, \textquotedblleft{}wide\textquotedblright{}, and \emph{ba}, \textquotedblleft{}place\textquotedblright{}, both inherited rather than borrowed (JMdict, KANJIDIC2).
\end{cousinweb}

\begin{grammarlens}[title={What you've built}]
One more word for everyday life at A2.
\end{grammarlens}

\subsection*{Your turn}
\begin{itemize}
\raggedright
  \item \emph{Say it:} \emph{hiroba}
  \item \emph{Say it:} \emph{hiroba}, once more
  \item \emph{Say it:} \emph{hiroba de matsuri ga arimasu} — There is a festival in the square.
  \item \emph{Recall:} say the Japanese for to go round, then the Japanese for city hall, then say \emph{hiroba} again
\end{itemize}

\begin{tcolorbox}[breakable,colback=teal!4,colframe=teal!35!black,title={Before you move on}]
What does \ja{ひろば} mean? (\textquotedblleft{}Town square; open space, plaza\textquotedblright{}.) Say it once more.
\end{tcolorbox}
\section[\texorpdfstring{\ja{こうじょう} (kōjō)}{こうじょう (kōjō)}]{\ja{こうじょう} (kōjō) — factory; plant, workshop}

\label{lesson:JA-C197-koujou}



\begin{quote}
Before the new one: say the Japanese for city hall, then the Japanese for town square.
\end{quote}

\subsection*{You'll want to know: \ja{こうじょう}}
\textbf{\ja{こうじょう}} — \emph{kōjō} — \textquotedblleft{}factory; plant, workshop\textquotedblright{}.

\begin{quote}
\textbf{\ja{ちちは} \ja{こうじょうで} \ja{はたらきます。}} — \emph{chichi wa kōjō de hatarakimasu} — \textquotedblleft{}My father works at a factory.\textquotedblright{}
\end{quote}

\begin{cousinweb}
\textbf{\ja{こうじょう}} (\emph{kōjō}) is Sino-Japanese (kango): two Chinese characters read with their on'yomi, readings borrowed from Chinese. They are \emph{kō}, \textquotedblleft{}craft, work\textquotedblright{} (Middle Chinese \emph{kuwng}) and \emph{jō}, \textquotedblleft{}place\textquotedblright{} (\emph{drjang}) (KANJIDIC2; Middle Chinese in Baxter's transcription).
\end{cousinweb}

\begin{grammarlens}[title={What you've built}]
One more word for everyday life at A2.
\end{grammarlens}

\subsection*{Your turn}
\begin{itemize}
\raggedright
  \item \emph{Say it:} \emph{kōjō}
  \item \emph{Say it:} \emph{kōjō}, once more
  \item \emph{Say it:} \emph{chichi wa kōjō de hatarakimasu} — My father works at a factory.
  \item \emph{Recall:} say the Japanese for city hall, then the Japanese for town square, then say \emph{kōjō} again
\end{itemize}

\begin{tcolorbox}[breakable,colback=teal!4,colframe=teal!35!black,title={Before you move on}]
What does \ja{こうじょう} mean? (\textquotedblleft{}Factory; plant, workshop\textquotedblright{}.) Say it once more.
\end{tcolorbox}
\section[\texorpdfstring{\ja{たてる} (tateru)}{たてる (tateru)}]{\ja{たてる} (tateru) — to build; to put up}

\label{lesson:JA-C197-tateru}



\begin{quote}
Before the new one: say the Japanese for town square, then the Japanese for factory.
\end{quote}

\subsection*{You'll want to know: \ja{たてる}}
\textbf{\ja{たてる}} — \emph{tateru} — \textquotedblleft{}to build; to put up\textquotedblright{}.

\begin{quote}
\textbf{\ja{えきの} \ja{まえに} \ja{ひゃっかてんを} \ja{たてます。}} — \emph{eki no mae ni hyakkaten o tatemasu} — \textquotedblleft{}They will build a department store in front of the station.\textquotedblright{}
\end{quote}

Polite \textbf{\ja{たてます}}, \emph{tatemasu}; an ichidan verb.

\begin{cousinweb}
\textbf{\ja{たてる}} (\emph{tateru}) is a native Japanese word (wago), inherited rather than borrowed from Chinese; its kanji was attached later as a kun'yomi (JMdict, KANJIDIC2).
\end{cousinweb}

\begin{grammarlens}[title={What you've built}]
One more word for everyday life at A2.
\end{grammarlens}

\subsection*{Your turn}
\begin{itemize}
\raggedright
  \item \emph{Say it:} \emph{tateru}
  \item \emph{Say it:} \emph{tateru}, once more
  \item \emph{Say it:} \emph{eki no mae ni hyakkaten o tatemasu} — They will build a department store in front of the station.
  \item \emph{Recall:} say the Japanese for town square, then the Japanese for factory, then say \emph{tateru} again
\end{itemize}

\begin{tcolorbox}[breakable,colback=teal!4,colframe=teal!35!black,title={Before you move on}]
What does \ja{たてる} mean? (\textquotedblleft{}To build; to put up\textquotedblright{}.) Say it once more.
\end{tcolorbox}
\section[\texorpdfstring{\ja{みつかる} (mitsukaru)}{みつかる (mitsukaru)}]{\ja{みつかる} (mitsukaru) — to be found; to turn up}

\label{lesson:JA-C197-mitsukaru}



\begin{quote}
Before the new one: say the Japanese for factory, then the Japanese for to build.
\end{quote}

\subsection*{You'll want to know: \ja{みつかる}}
\textbf{\ja{みつかる}} — \emph{mitsukaru} — \textquotedblleft{}to be found; to turn up\textquotedblright{}.

\begin{quote}
\textbf{\ja{こうばんで} \ja{さいふが} \ja{みつかりました。}} — \emph{kōban de saifu ga mitsukarimashita} — \textquotedblleft{}My wallet turned up at the police box.\textquotedblright{}
\end{quote}

Polite \textbf{\ja{みつかります}}, \emph{mitsukarimasu}; a godan verb. Its transitive partner \textbf{\ja{みつける}} is taught already.

\begin{cousinweb}
\textbf{\ja{みつかる}} (\emph{mitsukaru}) is a native Japanese word, inherited rather than borrowed, built on the same stem as \emph{mitsukeru}, \textquotedblleft{}to find\textquotedblright{} (JMdict, KANJIDIC2).
\end{cousinweb}

\begin{grammarlens}[title={What you've built}]
One more word for everyday life at A2.
\end{grammarlens}

\subsection*{Your turn}
\begin{itemize}
\raggedright
  \item \emph{Say it:} \emph{mitsukaru}
  \item \emph{Say it:} \emph{mitsukaru}, once more
  \item \emph{Say it:} \emph{kōban de saifu ga mitsukarimashita} — My wallet turned up at the police box.
  \item \emph{Recall:} say the Japanese for factory, then the Japanese for to build, then say \emph{mitsukaru} again
\end{itemize}

\begin{tcolorbox}[breakable,colback=teal!4,colframe=teal!35!black,title={Before you move on}]
What does \ja{みつかる} mean? (\textquotedblleft{}To be found; to turn up\textquotedblright{}.) Say it once more.
\end{tcolorbox}
\section[(dialogue)]{Town and street words, a review}

\label{lesson:JA-R197-review-196-197}



\begin{quote}
Say the Japanese for to head for and neighbourhood.
\end{quote}

\subsection*{Your turn}
Say these aloud:

\begin{itemize}
\raggedright
  \item to head for, neighbourhood, address, map, to go round
  \item city hall, town square, factory, to build, to be found
\end{itemize}

\begin{tcolorbox}[breakable,colback=teal!4,colframe=teal!35!black,title={Before you move on}]
To head for? (\textbf{\ja{むかう}}, \emph{mukau}.) Neighbourhood? (\textbf{\ja{きんじょ}}, \emph{kinjo}.) Address? (\textbf{\ja{じゅうしょ}}, \emph{jūsho}.) Map? (\textbf{\ja{ちず}}, \emph{chizu}.) To go round? (\textbf{\ja{まわる}}, \emph{mawaru}.) City hall? (\textbf{\ja{しやくしょ}}, \emph{shiyakusho}.) Town square? (\textbf{\ja{ひろば}}, \emph{hiroba}.) Factory? (\textbf{\ja{こうじょう}}, \emph{kōjō}.) To build? (\textbf{\ja{たてる}}, \emph{tateru}.) To be found? (\textbf{\ja{みつかる}}, \emph{mitsukaru}.)
\end{tcolorbox}
\section[(dialogue)]{City Hall, Squares and Factories — once more}

\label{lesson:JA-R197-review-final-197}



\begin{quote}
Say the Japanese for city hall and town square.
\end{quote}

\subsection*{Your turn}
\begin{itemize}
\raggedright
  \item \emph{Say it:} city hall, town square, factory, to build, to be found
\end{itemize}

\begin{tcolorbox}[breakable,colback=teal!4,colframe=teal!35!black,title={Before you move on}]
City hall? (\textbf{\ja{しやくしょ}}, \emph{shiyakusho}.) Town square? (\textbf{\ja{ひろば}}, \emph{hiroba}.) Factory? (\textbf{\ja{こうじょう}}, \emph{kōjō}.) To build? (\textbf{\ja{たてる}}, \emph{tateru}.) To be found? (\textbf{\ja{みつかる}}, \emph{mitsukaru}.)
\end{tcolorbox}
